\section{AHSME 1963}
        \begin{enumerate}
        	\item (\textit{AHSME 1963}) Which one of the following points is not on the graph of $y=\dfrac{x}{x+1}$?


 $\textbf{(A)}\ (0,0)\qquad \textbf{(B)}\ \left(-\frac{1}{2},-1\right)\qquad \textbf{(C)}\ \left(\frac{1}{2},\frac{1}{3}\right)\qquad \textbf{(D)}\ (-1,1)\qquad \textbf{(E)}\ (-2,2)$


 

	\item (\textit{AHSME 1963}) let $n=x-y^{x-y}$. Find $n$ when $x=2$ and $y=-2$.


 $\textbf{(A)}\ -14 \qquad \textbf{(B)}\ 0 \qquad \textbf{(C)}\ 1 \qquad \textbf{(D)}\ 18 \qquad \textbf{(E)}\ 256$


 

	\item (\textit{AHSME 1963}) If the reciprocal of $x+1$ is $x-1$, then $x$ equals:


 $\textbf{(A)}\ 0\qquad \textbf{(B)}\ 1\qquad \textbf{(C)}\ -1\qquad \textbf{(D)}\ \pm 1\qquad \textbf{(E)}\ \text{none of these}$


 

	\item (\textit{AHSME 1963}) For what value(s) of $k$ does the pair of equations $y=x^2$ and $y=3x+k$ have two identical solutions?


 $\textbf{(A)}\ \frac{4}{9}\qquad \textbf{(B)}\ -\frac{4}{9}\qquad \textbf{(C)}\ \frac{9}{4}\qquad \textbf{(D)}\ -\frac{9}{4}\qquad \textbf{(E)}\ \pm\frac{9}{4}$


 

	\item (\textit{AHSME 1963}) If $x$ and $\log_{10} x$ are real numbers and $\log_{10} x<0$, then:


 $\textbf{(A)}\ x<0 \qquad \textbf{(B)}\ -1<x<1 \qquad \textbf{(C)}\ 0<x\le 1 \\ \textbf{(D)}\ -1<x<0 \qquad \textbf{(E)}\ 0<x<1$


 

	\item (\textit{AHSME 1963}) $\triangle BAD$ is right-angled at $B$. On $AD$ there is a point $C$ for which $AC=CD$ and $AB=BC$. The magnitude of $\angle DAB$ is:


 $\textbf{(A)}\ 67\frac{1}{2}^{\circ}\qquad \textbf{(B)}\ 60^{\circ}\qquad \textbf{(C)}\ 45^{\circ}\qquad \textbf{(D)}\ 30^{\circ}\qquad \textbf{(E)}\ 22\frac{1}{2}^{\circ}$


 

	\item (\textit{AHSME 1963}) Given the four equations:


 $\textbf{(1)}\ 3y-2x=12 \qquad\textbf{(2)}\ -2x-3y=10 \qquad\textbf{(3)}\ 3y+2x=12 \qquad\textbf{(4)}\ 2y+3x=10$


 The pair representing the perpendicular lines is:


 

$\textbf{(A)}\ \text{(1) and (4)}\qquad \textbf{(B)}\ \text{(1) and (3)}\qquad \textbf{(C)}\ \text{(1) and (2)}\qquad \textbf{(D)}\ \text{(2) and (4)}\qquad \textbf{(E)}\ \text{(2) and (3)}$


 

	\item (\textit{AHSME 1963}) The smallest positive integer $x$ for which $1260x=N^3$, where $N$ is an integer, is:


 

$\textbf{(A)}\ 1050 \qquad \textbf{(B)}\ 1260 \qquad \textbf{(C)}\ 1260^2 \qquad \textbf{(D)}\ 7350 \qquad \textbf{(E)}\ 44100$


 

	\item (\textit{AHSME 1963}) In the expansion of $\left(a-\dfrac{1}{\sqrt{a}}\right)^7$ the coefficient of $a^{-\dfrac{1}{2}}$ is:


 $\textbf{(A)}\ -7 \qquad \textbf{(B)}\ 7 \qquad \textbf{(C)}\ -21 \qquad \textbf{(D)}\ 21 \qquad \textbf{(E)}\ 35$


 

	\item (\textit{AHSME 1963}) Point $P$ is taken interior to a square with side-length $a$ and such that is it equally distant from two 
consecutive vertices and from the side opposite these vertices. If $d$ represents the common distance, then $d$ equals:


 $\textbf{(A)}\ \frac{3a}{5}\qquad \textbf{(B)}\ \frac{5a}{8}\qquad \textbf{(C)}\ \frac{3a}{8}\qquad \textbf{(D)}\ \frac{a\sqrt{2}}{2}\qquad \textbf{(E)}\ \frac{a}{2}$


 

	\item (\textit{AHSME 1963}) The arithmetic mean of a set of $50$ numbers is $38$. If two numbers of the set, namely $45$ and $55$, are discarded, 
the arithmetic mean of the remaining set of numbers is:


 $\textbf{(A)}\ 38.5 \qquad \textbf{(B)}\ 37.5 \qquad \textbf{(C)}\ 37 \qquad \textbf{(D)}\ 36.5 \qquad \textbf{(E)}\ 36$


 

	\item (\textit{AHSME 1963}) Three vertices of parallelogram $PQRS$ are $P(-3,-2), Q(1,-5), R(9,1)$ with $P$ and $R$ diagonally opposite. 
The sum of the coordinates of vertex $S$ is:


 $\textbf{(A)}\ 13 \qquad \textbf{(B)}\ 12 \qquad \textbf{(C)}\ 11 \qquad \textbf{(D)}\ 10 \qquad \textbf{(E)}\ 9$


 

	\item (\textit{AHSME 1963}) If $2^a+2^b=3^c+3^d$, the number of integers $a,b,c,d$ which can possibly be negative, is, at most:


 $\textbf{(A)}\ 4 \qquad \textbf{(B)}\ 3 \qquad \textbf{(C)}\ 2 \qquad \textbf{(D)}\ 1 \qquad \textbf{(E)}\ 0$


 

	\item (\textit{AHSME 1963}) Given the equations $x^2+kx+6=0$ and $x^2-kx+6=0$. If, when the roots of the equation are suitably listed, 
each root of the second equation is $5$ more than the corresponding root of the first equation, then $k$ equals:


 $\textbf{(A)}\ 5 \qquad \textbf{(B)}\ -5 \qquad \textbf{(C)}\ 7 \qquad \textbf{(D)}\ -7 \qquad \textbf{(E)}\ \text{none of these}$


 

	\item (\textit{AHSME 1963}) A circle is inscribed in an equilateral triangle, and a square is inscribed in the circle. 
The ratio of the area of the triangle to the area of the square is:


 $\textbf{(A)}\ \sqrt{3}:1\qquad \textbf{(B)}\ \sqrt{3}:\sqrt{2}\qquad \textbf{(C)}\ 3\sqrt{3}:2\qquad \textbf{(D)}\ 3:\sqrt{2}\qquad \textbf{(E)}\ 3:2\sqrt{2}$


 

	\item (\textit{AHSME 1963}) Three numbers $a,b,c$, none zero, form an arithmetic progression. Increasing $a$ by $1$ or increasing $c$ by $2$ results 
in a geometric progression. Then $b$ equals:


 $\textbf{(A)}\ 16 \qquad \textbf{(B)}\ 14 \qquad \textbf{(C)}\ 12 \qquad \textbf{(D)}\ 10 \qquad \textbf{(E)}\ 8$


 

	\item (\textit{AHSME 1963}) The expression $\dfrac{\dfrac{a}{a+y}+\dfrac{y}{a-y}}{\dfrac{y}{a+y}-\dfrac{a}{a-y}}$, $a$ real, $a\neq 0$, has the value $-1$ for:


 $\textbf{(A)}\ \text{all but two real values of }y \qquad \\ \textbf{(B)}\ \text{only two real values of }y \qquad \\ \textbf{(C)}\ \text{all real values of }y\qquad \\ \textbf{(D)}\ \text{only one real value of }y\qquad \\ \textbf{(E)}\ \text{no real values of }y$


 

	\item (\textit{AHSME 1963}) Chord $EF$ is the perpendicular bisector of chord $BC$, intersecting it in $M$. 
Between $B$ and $M$ point $U$ is taken, 
and $EU$ extended meets the circle in $A$. Then, for any selection of $U$, as described, $\triangle EUM$ is similar to:


 
                    \begin{figure}[H]
                        \centering
                        \includegraphics[width=0.5\textwidth]{Problem/AHSME_I_1963_18_Problem_Figure_1.png}
                    \end{figure}
                    


 

$\textbf{(A)}\ \triangle EFA \qquad \textbf{(B)}\ \triangle EFC \qquad \textbf{(C)}\ \triangle ABM \qquad \textbf{(D)}\ \triangle ABU \qquad \textbf{(E)}\ \triangle FMC$


 

	\item (\textit{AHSME 1963}) In counting $n$ colored balls, some red and some black, it was found that $49$ of the first $50$ counted were red. 
Thereafter, $7$ out of every $8$ counted were red. If, in all, $90$ \% or more of the balls counted were red, the maximum value of $n$ is:


 $\textbf{(A)}\ 225 \qquad \textbf{(B)}\ 210 \qquad \textbf{(C)}\ 200 \qquad \textbf{(D)}\ 180 \qquad \textbf{(E)}\ 175$


 

	\item (\textit{AHSME 1963}) Two men at points $R$ and $S$, $76$ miles apart, set out at the same time to walk towards each other. 
The man at $R$ walks uniformly at the rate of $4\frac{1}{2}$ miles per hour; the man at $S$ walks at the constant 
rate of $3\frac{1}{4}$ miles per hour for the first hour, at $3\frac{3}{4}$ miles per hour for the second hour,
and so on, in arithmetic progression. If the men meet $x$ miles nearer $R$ than $S$ in an integral number of hours, then $x$ is:


 $\textbf{(A)}\ 10 \qquad \textbf{(B)}\ 8 \qquad \textbf{(C)}\ 6 \qquad \textbf{(D)}\ 4 \qquad \textbf{(E)}\ 2$


 

	\item (\textit{AHSME 1963}) The expression $x^2-y^2-z^2+2yz+x+y-z$ has:


 $\textbf{(A)}\ \text{no linear factor with integer coefficients and integer exponents} \qquad \\ \textbf{(B)}\ \text{the factor }-x+y+z \qquad \\ \textbf{(C)}\ \text{the factor }x-y-z+1 \qquad \\ \textbf{(D)}\ \text{the factor }x+y-z+1 \qquad \\ \textbf{(E)}\ \text{the factor }x-y+z+1$


 

	\item (\textit{AHSME 1963}) Acute-angled $\triangle ABC$ is inscribed in a circle with center at $O$; $\stackrel \frown {AB} = 120$ and $\stackrel \frown {BC} = 72$. 
$A$ point $E$ is taken in minor arc $AC$ such that $OE$ is perpendicular to $AC$. Then the ratio of the magnitudes of $\angle OBE$ and $\angle BAC$ is:


 $\textbf{(A)}\ \frac{5}{18}\qquad \textbf{(B)}\ \frac{2}{9}\qquad \textbf{(C)}\ \frac{1}{4}\qquad \textbf{(D)}\ \frac{1}{3}\qquad \textbf{(E)}\ \frac{4}{9}$


 

	\item (\textit{AHSME 1963}) A gives $B$ as many cents as $B$ has and $C$ as many cents as $C$ has. Similarly, $B$ then gives $A$ and $C$ as many cents as each then has. 
$C$, similarly, then gives $A$ and $B$ as many cents as each then has. If each finally has $16$ cents, with how many cents does $A$ start?


 $\textbf{(A)}\ 24 \qquad \textbf{(B)}\ 26\qquad \textbf{(C)}\ 28 \qquad \textbf{(D)}\ 30 \qquad \textbf{(E)}\ 32$


 

	\item (\textit{AHSME 1963}) Consider equations of the form $x^2 + bx + c = 0$. How many such equations have real roots and have coefficients $b$ and $c$ selected 
from the set of integers $\{1,2,3, 4, 5,6\}$?


 $\textbf{(A)}\ 20 \qquad \textbf{(B)}\ 19 \qquad \textbf{(C)}\ 18 \qquad \textbf{(D)}\ 17 \qquad \textbf{(E)}\ 16$


 

	\item (\textit{AHSME 1963}) Point $F$ is taken in side $AD$ of square $ABCD$. At $C$ a perpendicular is drawn to $CF$, meeting $AB$ extended at $E$. 
The area of $ABCD$ is $256$ square inches and the area of $\triangle CEF$ is $200$ square inches. Then the number of inches in $BE$ is:


 
                    \begin{figure}[H]
                        \centering
                        \includegraphics[width=0.5\textwidth]{Problem/AHSME_I_1963_25_Problem_Figure_1.png}
                    \end{figure}
                    


 $\textbf{(A)}\ 12 \qquad \textbf{(B)}\ 14 \qquad \textbf{(C)}\ 15 \qquad \textbf{(D)}\ 16 \qquad \textbf{(E)}\ 20$


 

	\item (\textit{AHSME 1963}) Version 1
Consider the statements:


 $\textbf{(1)}\ p\text{ }\wedge\sim q\wedge r\qquad\textbf{(2)}\ \sim p\text{ }\wedge\sim q\wedge r\qquad\textbf{(3)}\ p\text{ }\wedge\sim q\text{ }\wedge\sim r\qquad\textbf{(4)}\ \sim p\text{ }\wedge q\text{ }\wedge r$


 where $p,q$, and $r$ are propositions. How many of these imply the truth of $(p\rightarrow q)\rightarrow r$?


 $\textbf{(A)}\ 0 \qquad \textbf{(B)}\ 1\qquad \textbf{(C)}\ 2 \qquad \textbf{(D)}\ 3 \qquad \textbf{(E)}\ 4$


 Version 2
Consider the statements (1) $p$ and $r$ are true and $q$ is false (2) $r$ is true and $p$ and $q$ are false (3) $p$ is true and $q$ and $r$ are false (4) $q$ and $r$ are true and $p$ is false. 
How many of these imply the truth of the statement
"$r$ is implied by the statement that $p$ implies $q$"?


 $\textbf{(A)}\ 0 \qquad \textbf{(B)}\ 1\qquad \textbf{(C)}\ 2 \qquad \textbf{(D)}\ 3 \qquad \textbf{(E)}\ 4$


 

	\item (\textit{AHSME 1963}) Six straight lines are drawn in a plane with no two parallel and no three concurrent. The number of regions into which they divide the plane is:


 $\textbf{(A)}\ 16 \qquad \textbf{(B)}\ 20\qquad \textbf{(C)}\ 22 \qquad \textbf{(D)}\ 24 \qquad \textbf{(E)}\ 26$


 

	\item (\textit{AHSME 1963}) Given the equation $3x^2 - 4x + k = 0$ with real roots. The value of $k$ for which the product of the roots of the equation is a maximum is:


 $\textbf{(A)}\ \frac{16}{9}\qquad \textbf{(B)}\ \frac{16}{3}\qquad \textbf{(C)}\ \frac{4}{9}\qquad \textbf{(D)}\ \frac{4}{3}\qquad \textbf{(E)}\ -\frac{4}{3}$


 

	\item (\textit{AHSME 1963}) A particle projected vertically upward reaches, at the end of $t$ seconds, an elevation of $s$ feet where $s = 160 t - 16t^2$. The highest elevation is:


 $\textbf{(A)}\ 800 \qquad \textbf{(B)}\ 640\qquad \textbf{(C)}\ 400 \qquad \textbf{(D)}\ 320 \qquad \textbf{(E)}\ 160$


 

	\item (\textit{AHSME 1963}) Let $F=\log\dfrac{1+x}{1-x}$. Find a new function $G$ by replacing each $x$ in $F$ by $\dfrac{3x+x^3}{1+3x^2}$, and simplify. 
The simplified expression $G$ is equal to:


 $\textbf{(A)}\ -F \qquad \textbf{(B)}\ F\qquad \textbf{(C)}\ 3F \qquad \textbf{(D)}\ F^3 \qquad \textbf{(E)}\ F^3-F$


 

	\item (\textit{AHSME 1963}) The number of solutions in positive integers of $2x+3y=763$ is:


 $\textbf{(A)}\ 255 \qquad \textbf{(B)}\ 254\qquad \textbf{(C)}\ 128 \qquad \textbf{(D)}\ 127 \qquad \textbf{(E)}\ 0$


 

	\item (\textit{AHSME 1963}) The dimensions of a rectangle $R$ are $a$ and $b$, $a < b$. It is required to obtain a rectangle with dimensions $x$ and $y$, $x < a, y < a$, 
so that its perimeter is one-third that of $R$, and its area is one-third that of $R$. The number of such (different) rectangles is:


 $\textbf{(A)}\ 0 \qquad \textbf{(B)}\ 1\qquad \textbf{(C)}\ 2 \qquad \textbf{(D)}\ 4 \qquad \textbf{(E)}\ \infty$


 

	\item (\textit{AHSME 1963}) Given the line $y = \dfrac{3}{4}x + 6$ and a line $L$ parallel to the given line and $4$ units from it. A possible equation for $L$ is:


 $\textbf{(A)}\ y =\frac{3}{4}x+1\qquad \textbf{(B)}\ y =\frac{3}{4}x\qquad \textbf{(C)}\ y =\frac{3}{4}x-\frac{2}{3}\qquad  \\ \textbf{(D)}\ y = \dfrac{3}{4}x -1 \qquad \textbf{(E)}\ y = \dfrac{3}{4}x + 2$


 

	\item (\textit{AHSME 1963}) In $\triangle ABC$, side $a = \sqrt{3}$, side $b = \sqrt{3}$, and side $c > 3$. Let $x$ be the largest number such that the magnitude, 
in degrees, of the angle opposite side $c$ exceeds $x$. Then $x$ equals:


 $\textbf{(A)}\ 150^{\circ} \qquad \textbf{(B)}\ 120^{\circ}\qquad \textbf{(C)}\ 105^{\circ} \qquad \textbf{(D)}\ 90^{\circ} \qquad \textbf{(E)}\ 60^{\circ}$


 

	\item (\textit{AHSME 1963}) The lengths of the sides of a triangle are integers, and its area is also an integer. 
One side is $21$ and the perimeter is $48$. The shortest side is:


 $\textbf{(A)}\ 8 \qquad \textbf{(B)}\ 10\qquad \textbf{(C)}\ 12 \qquad \textbf{(D)}\ 14 \qquad \textbf{(E)}\ 16$


 

	\item (\textit{AHSME 1963}) A person starting with $\$64$ and making $6$ bets, wins three times and loses three times, 
the wins and losses occurring in random order. The chance for a win is equal to the chance for a loss. 
If each wager is for half the money remaining at the time of the bet, then the final result is:


 $\textbf{(A)}\text{ a loss of }\$ 27 \qquad \textbf{(B)}\text{ a gain of }\$ 27 \qquad \textbf{(C)}\text{ a loss of }\$ 37 \qquad \\ \textbf{(D)}\text{ neither a gain nor a loss}\qquad \\ \textbf{(E)}\text{ a gain or a loss depending upon the order in which the wins and losses occur}$


 

	\item (\textit{AHSME 1963}) Given points $P_1, P_2,\cdots,P_7$ on a straight line, in the order stated (not necessarily evenly spaced). 
Let $P$ be an arbitrarily selected point on the line and let $s$ be the sum of the undirected lengths
$PP_1, PP_2, \cdots , PP_7$. Then $s$ is smallest if and only if the point $P$ is:


 $\textbf{(A)}\ \text{midway between }P_1\text{ and }P_7\qquad \\ \textbf{(B)}\ \text{midway between }P_2\text{ and }P_6\qquad \\ \textbf{(C)}\ \text{midway between }P_3\text{ and }P_5\qquad \\ \textbf{(D)}\ \text{at }P_4 \qquad \textbf{(E)}\ \text{at }P_1$


 

	\item (\textit{AHSME 1963}) Point $F$ is taken on the extension of side $AD$ of parallelogram $ABCD$. $BF$ intersects diagonal $AC$ at $E$ and side $DC$ at $G$. 
If $EF = 32$ and $GF = 24$, then $BE$ equals:


 
                    \begin{figure}[H]
                        \centering
                        \includegraphics[width=0.5\textwidth]{Problem/AHSME_I_1963_38_Problem_Figure_1.png}
                    \end{figure}
                    


 $\textbf{(A)}\ 4 \qquad \textbf{(B)}\ 8\qquad \textbf{(C)}\ 10 \qquad \textbf{(D)}\ 12 \qquad \textbf{(E)}\ 16$


 

	\item (\textit{AHSME 1963}) In $\triangle ABC$ lines $CE$ and $AD$ are drawn so that 
$\dfrac{CD}{DB}=\dfrac{3}{1}$ and $\dfrac{AE}{EB}=\dfrac{3}{2}$. Let $r=\dfrac{CP}{PE}$ 
where $P$ is the intersection point of $CE$ and $AD$. Then $r$ equals:


 
                    \begin{figure}[H]
                        \centering
                        \includegraphics[width=0.5\textwidth]{Problem/AHSME_I_1963_39_Problem_Figure_1.png}
                    \end{figure}
                    


 

$\textbf{(A)}\ 3 \qquad \textbf{(B)}\ \dfrac{3}{2}\qquad \textbf{(C)}\ 4 \qquad \textbf{(D)}\ 5 \qquad \textbf{(E)}\ \dfrac{5}{2}$


 

	\item (\textit{AHSME 1963}) If $x$ is a number satisfying the equation $\sqrt[3]{x+9}-\sqrt[3]{x-9}=3$, then $x^2$ is between:


 $\textbf{(A)}\ 55\text{ and }65\qquad \textbf{(B)}\ 65\text{ and }75\qquad \textbf{(C)}\ 75\text{ and }85\qquad \textbf{(D)}\ 85\text{ and }95\qquad \textbf{(E)}\ 95\text{ and }105$


 

\end{enumerate}